\documentclass[11pt,a4paper]{article}
\usepackage[T1]{fontenc}
\usepackage{lmodern}
\usepackage{amsmath,amssymb,booktabs,geometry,hyperref}
\geometry{margin=25mm}
\hypersetup{hidelinks}
\title{chem-kumo: Mathematical Foundations and Numerical Algorithms}
\author{Matteo Tommasini}
\date{Version 1.2.0-dev}
\begin{document}
\maketitle
\tableofcontents

\section{Scope and conventions}
KUMO is a non-SCF evaluator of real ground-state electronic fields from
Gaussian formatted checkpoint (FCHK) data. It neither determines the electronic
state nor relaxes orbital coefficients. All derivatives below are with respect
to the electronic Cartesian coordinate $\mathbf r$, with the molecular geometry
and FCHK data fixed. No nuclear derivatives or electronic response equations
are implemented. Atomic units are used: distances in bohr, energies in hartree.
AO and MO values have dimensions $a_0^{-3/2}$, density $a_0^{-3}$,
its gradient $a_0^{-4}$, Hessian and Laplacian $a_0^{-5}$, and the positive
kinetic-energy density $E_h a_0^{-3}$ in atomic units.

\section{Gaussian atomic-orbital basis}
A Cartesian primitive centered at $\mathbf A$ is
\begin{equation}
 g_{\ell mn}(\mathbf r;\alpha)=N_{\ell mn}(\alpha)
 u^\ell v^m w^n\exp[-\alpha(u^2+v^2+w^2)],\qquad
 (u,v,w)=\mathbf r-\mathbf A.
\end{equation}
An AO is a contracted linear combination of primitives (with the appropriate
normalization and contraction conventions):
\begin{equation}
 \chi_\mu(\mathbf r)=\sum_p d_{\mu p}g_{\mu p}(\mathbf r).
\end{equation}
For pure spherical shells, a fixed Cartesian-to-real-spherical linear
transformation maps primitive or contracted Cartesian components into the
FCHK AO order. This same transformation is applied to all spatial derivatives.
For combined SP shells, S and P contractions use their respective coefficients.
The FCHK record for SP contraction coefficients need not exist when no SP
shell is present. The AO ordering and normalization must agree with the
coefficient and density matrices read from the same FCHK.

\section{Analytic AO derivatives through second order}
Define the polynomial factors
\begin{align}
 f_\ell(u)&=u^\ell,\\
 D_\ell(u;\alpha)&=\ell u^{\ell-1}-2\alpha u^{\ell+1},\\
 E_\ell(u;\alpha)&=\ell(\ell-1)u^{\ell-2}
 -2\alpha(2\ell+1)u^\ell+4\alpha^2u^{\ell+2}.
\end{align}
Terms whose prefactors vanish are evaluated as zero, without forming
negative powers at the origin. With $G=N_{\ell mn}\exp[-\alpha(u^2+v^2+w^2)]$,
\begin{align}
 \partial_x g_{\ell mn}&=G D_\ell(u;\alpha)f_m(v)f_n(w),\\
 \partial_{xx}g_{\ell mn}&=G E_\ell(u;\alpha)f_m(v)f_n(w),\\
 \partial_{xy}g_{\ell mn}&=G D_\ell(u;\alpha)D_m(v;\alpha)f_n(w).
\end{align}
The other Cartesian components follow by permutation. Derivation commutes
with contraction and the Cartesian-to-pure transformation. The AO evaluator
returns $\chi_\mu$, three first derivatives, and six unique second derivatives
in the order $(xx,yy,zz,xy,xz,yz)$. This is the storage convention, not the
matrix layout of the symmetric $3\times3$ Hessian.

\section{Molecular orbitals}
For real molecular orbitals in spin channel $\sigma\in\{\alpha,\beta\}$,
\begin{align}
 \psi_i^\sigma(\mathbf r)&=\sum_\mu C_{\mu i}^\sigma\chi_\mu(\mathbf r),\\
 \partial_a\psi_i^\sigma&=\sum_\mu C_{\mu i}^\sigma\partial_a\chi_\mu,\\
 \partial_{ab}\psi_i^\sigma&=\sum_\mu C_{\mu i}^\sigma\partial_{ab}\chi_\mu.
\end{align}
These are derivatives of the signed orbital amplitude, not of its square.
The public MO indices are one-based, whereas internal coefficient matrices
use zero-based indexing. Alpha and beta select orbital channels, not total
or spin-difference orbitals. No MO phase alignment is applied.

\section{Spin densities and the AO density matrix}
Let $\boldsymbol\chi(\mathbf r)$ be the column vector of AO values and
$P^\sigma$ a real symmetric spin density matrix. Then
\begin{equation}
 \rho_\sigma=\boldsymbol\chi^T P^\sigma\boldsymbol\chi
 =\sum_i n_{i\sigma}(\psi_i^\sigma)^2,
 \qquad P^\sigma=C^\sigma n^\sigma(C^\sigma)^T.
\end{equation}
Here $n^\sigma$ is diagonal and contains spin-orbital occupations. In the
restricted closed-shell case the two spin matrices are equal and each
contains half the total density. Define
\begin{equation}
 P^{\mathrm{total}}=P^\alpha+P^\beta,\qquad
 P^{\mathrm{spin}}=P^\alpha-P^\beta.
\end{equation}
KUMO supports total, alpha, beta, and alpha-minus-beta components.
The spin-difference field is signed and need not be nonnegative.

\section{Spatial derivatives of the electron density}
Write $\boldsymbol\chi_a=\partial_a\boldsymbol\chi$ and
$\boldsymbol\chi_{ab}=\partial_a\partial_b\boldsymbol\chi$.
For any of the symmetric matrices $P^s$ described above,
\begin{align}
 \rho_s&=\boldsymbol\chi^T P^s\boldsymbol\chi,\\
 \partial_a\rho_s&=2\boldsymbol\chi_a^T P^s\boldsymbol\chi,\\
 \partial_{ab}\rho_s&=2\boldsymbol\chi_{ab}^T P^s\boldsymbol\chi
       +2\boldsymbol\chi_a^T P^s\boldsymbol\chi_b,\\
 \nabla^2\rho_s&=\sum_{a=x,y,z}\partial_{aa}\rho_s
                  =\operatorname{tr}H(\rho_s).
\end{align}
The full Hessian is symmetric; the off-diagonal components are evaluated
analytically. The positive-definite kinetic-energy density for a physical
spin channel is
\begin{equation}
 \tau_\sigma(\mathbf r)=\frac12\sum_{a=x,y,z}
 \boldsymbol\chi_a^T P^\sigma\boldsymbol\chi_a
 =\frac12\sum_i n_{i\sigma}|\nabla\psi_i^\sigma|^2.
\end{equation}
The same linear formula applied to $P^{\mathrm{spin}}$ yields
$\tau_\alpha-\tau_\beta$, which is not necessarily positive.
The identity
\begin{equation}
 \nabla^2\rho_s=2(\nabla^2\boldsymbol\chi)^TP^s\boldsymbol\chi
 +4\tau_s
\end{equation}
provides an additional algebraic consistency check.

\section{Blocked dense contractions with Eigen}
For $B$ spatial points, let $\Phi,\Phi_x,\Phi_y,\Phi_z$ be
$N_{\mathrm{AO}}\times B$ panels of AO values and first derivatives.
The implementation packs these into one column-major matrix
\begin{equation}
 A=[\Phi\ \Phi_x\ \Phi_y\ \Phi_z]
 \in\mathbb R^{N_{\mathrm{AO}}\times4B}
\end{equation}
and applies the selected density matrix in one Eigen matrix product:
\begin{equation}
 V=P^s A=[V_0\ V_x\ V_y\ V_z].
\end{equation}
For each point (column) $j$, the fields are obtained by vector products:
\begin{align}
 \rho_{s,j}&=\Phi_j^T(V_0)_j,\\
 (\nabla\rho_s)_{a,j}&=2(\Phi_a)_j^T(V_0)_j,\\
 (H_s)_{ab,j}&=2(\Phi_{ab})_j^T(V_0)_j
                  +2(\Phi_a)_j^T(V_b)_j,\\
 \tau_{s,j}&=\tfrac12\sum_a(\Phi_a)_j^T(V_a)_j.
\end{align}
Six second-derivative panels are retained for the linear reductions but
\emph{are not} multiplied by $P^s$. Only the diagonal pointwise contractions
are required; no $B\times B$ field matrix is formed. The AO derivatives
are still evaluated point by point inside each block.

The dense matrix product requires $O(4N_{\mathrm{AO}}^2B)$ arithmetic
per block; AO derivative evaluation and reductions scale approximately
linearly with $N_{\mathrm{AO}}$ for fixed basis structure. The fused product
reduces separate Eigen calls and may improve reuse of $P^s$ in cache,
but it does not change the asymptotic complexity or reduce the number of
mathematical density-matrix applications. No sparse screening is applied.

For a full block, ten AO panels (one value, three gradients, six Hessians)
and four product panels together require approximately
$14N_{\mathrm{AO}}B$ doubles, excluding the density matrices, output and
temporary storage. At $N_{\mathrm{AO}}=1410$, $B=128$, this is approximately
20.2 MB of panel data. The public blocked API returns a vector of results,
so its total output memory grows with the number of points even though AO
scratch data are block-local.

\section{Validation and benchmark methodology}
AO, MO, and density spatial derivatives are checked against centered finite
differences of independent value/gradient evaluations. Finite differences
are used only for testing, never to compute production derivatives.
Density-matrix results are also checked against independent occupied-MO
sums for density, gradient, Hessian, and $\tau$, and blocked results are
compared with the pointwise reference. The regression suite includes Gaussian
09 reference Cube comparisons for density and signed orbital amplitudes.
These tests establish agreement on the tested FCHK cases, not universal
support for every Gaussian basis or method.

The optional \texttt{kumo-spatial-benchmark} executable times AO values,
AO derivatives, representative dense contractions, density-only evaluation,
and pointwise/blocked density derivatives. The timer uses one warm-up and
the minimum of repeated wall-clock times. Representative contraction timings
are not additive components of end-to-end timings; the full methods have
different setup and allocation costs. Benchmark times are performance
measurements, never pass/fail correctness thresholds.

\section{Limits and intended use}
KUMO evaluates the real electronic fields defined by the supplied FCHK.
It does not perform SCF, orbital optimization, response theory, or nuclear
coordinate differentiation. The present blocked derivative evaluator uses
dense matrices and fixed-size blocks, without AO screening or streaming
output of derivative fields. Its mathematical derivatives are analytic
with respect to the supplied finite Gaussian representation; grid sampling
and floating-point arithmetic introduce their usual numerical limitations.
\end{document}
